\subsection{The reduced Ricci-flow test}\label{sec:cotton-variation}

At a smooth constant-scalar Bach-flat metric let \(g(t)\) start at
\(g\) under Ricci flow. Use the ordered-index tensor
\(C_{ijk}=\nabla_i Z_{jk}-\nabla_j Z_{ik}\), which is twice the
Cotton tensor formed from \((\Ric-Rg/6)/2\) with the same index order.
Write \(s_m=\operatorname{tr}((Z^\sharp)^m)\). The tangent to the
volume-normalized flow is \(-2Z\). The unnormalized tangent differs
by a pure trace, which is in the Hessian kernel at the critical
metric. The second derivative for either flow in its own time
variable is therefore \(4H_g(Z,Z)\). It reduces to
\begin{equation}\label{eq:ricci-cotton-hess}
 \left.\frac{d^2}{dt^2}E_+(g(t))\right|_{t=0}
 =\int_M\left[\frac7{12}s_2^2-s_4
      -3Z^{ia}C_{ijk}C_a{}^{jk}\right]dV_g .
\end{equation}
The quartic term is nonnegative. In a Ricci eigenframe with
\(Z^\sharp\) eigenvalues \(\xi_1,\ldots,\xi_4\), it is
\begin{equation}\label{eq:ricci-quartic}
 \frac7{12}s_2^2-s_4
 =\frac16\sum_{ij|kl}(\xi_i-\xi_j)^2(\xi_k-\xi_l)^2 ,
\end{equation}
where the sum runs over the three partitions into pairs. In fact,
if \(p=\xi_1+\xi_2\), \(q=\xi_1+\xi_3\) and \(r=\xi_1+\xi_4\),
then \(s_2=p^2+q^2+r^2\) and
\(s_4=3s_2^2/4-(p^4+q^4+r^4)/2\). Six times the quartic is
\((p^2-q^2)^2+(q^2-r^2)^2+(r^2-p^2)^2\). It vanishes exactly when
at least three eigenvalues coincide.

To derive \eqref{eq:ricci-cotton-hess}, the constant-scalar Bach
equation and \(\operatorname{div}Z=0\) give
\begin{equation}\label{eq:cotton-div-bach}
 \nabla^i C_{ijk}=-(W\cdot Z)_{jk} .
\end{equation}
Indeed, \(\nabla^i\nabla_j Z_{ik}=(\Ric\circ Z)_{jk}-(\Rm\cdot Z)_{jk}\),
while \(\Rm\cdot Z=W\cdot Z-Z^2-RZ/12+s_2g/2\), with \(Z^2=Z\circ Z^\sharp\).
Subtract the commuted derivative from the stationary formula for
\(\Delta Z\). The Weyl contraction \(\langle W\cdot Z,Z^2\rangle\)
vanishes pointwise in dimension four. To see this, set
\(w_{ij}=W_{ijij}\) in the eigenframe. Weyl trace-freeness gives
\(w_{ij}=w_{kl}\) for complementary pairs and
\(w_{12}+w_{13}+w_{14}=0\). The coefficient of each \(w_{ij}\) in
the contraction is
\[
 \xi_i\xi_j(\xi_i+\xi_j)+\xi_k\xi_l(\xi_k+\xi_l)
 =-\sum_{a<b<c}\xi_a\xi_b\xi_c .
\]
Integration by parts in \eqref{eq:cotton-div-bach} therefore gives
\(\int_M C_{ijk}\nabla_i(Z^2)_{jk}\,dV_g=0\).

For \(T_{ijk}=\nabla_iZ_{jk}\), let \(P_1,P_2,X_1,X_2\) denote
the sums of \(\xi_i T_{ijk}^2\), \(\xi_j T_{ijk}^2\),
\(\xi_i T_{ijk}T_{jik}\), and \(\xi_k T_{ijk}T_{jik}\), respectively.
Let \(J=Z^{ia}C_{ijk}C_a{}^{jk}=P_1+P_2-2X_1\).
The two derivative terms of \eqref{eq:ricci-hess} add to
\(5P_2-3J-2(X_1+X_2)\). The vanishing integrated pairing just
proved is \(\int_M(2P_2-X_1-X_2)\,dV_g=0\). Pairing the stationary
Bach equation with \(Z^2\) gives
\[
 2\int_M\Rm(Z,Z^2)\,dV_g
  =\int_M(2P_2+Rs_3/6+s_2^2/2)\,dV_g .
\]
The Weyl cancellation and the curvature decomposition give
\(\Rm(Z,Z^2)=-s_4-Rs_3/12+s_2^2/2\). Thus the integrated
derivative terms of \eqref{eq:ricci-hess} reduce to the integral
of \(-s_4-Rs_3/6+s_2^2/4-3J\). Its scalar terms then produce
\eqref{eq:ricci-cotton-hess}.

At an attained minimum, three times the integrated Cotton
contraction is at most the integrated quartic in
\eqref{eq:ricci-quartic}. A strict reverse inequality at a
stationary positive metric gives a same-type positive competitor
along short forward flow. The first derivative vanishes, and the
positive Yamabe sign persists by continuity of the first conformal
Laplacian eigenvalue. The contraction can have either sign at
the level of a first jet. For example, entries
\(T_{123}=T_{132}=T_{312}=T_{321}=h\) with all other entries zero
give \(J=(\xi_2-\xi_4)h^2\) and satisfy the trace and divergence
constraints on \(T\). They supply no control of the integrated
sign on an actual stationary metric in the surviving band.

\begin{proposition}[Cotton-flat positive constant-scalar gauge]
\label{prop:cotton-flat-bach}
Let \(g\) be a smooth Bach-flat Riemannian metric on a connected
four-manifold without boundary, of constant scalar curvature
\(R>0\). If \(C=0\), then \(g\) is Einstein or locally
conformally flat. On a closed smooth homotopy four-sphere the
locally conformally flat case is round up to scale and
diffeomorphism.
\end{proposition}

\begin{proof}
Equation~\eqref{eq:cotton-div-bach} gives \(W\cdot Z=0\).
Adding its diagonal components with indices 1 and 2 gives
\(p w_{12}=0\), using the same \(p,q,r\) and \(w_{ij}\) as above.
The other complementary pairs give \(q w_{13}=r w_{14}=0\).
If \(Z\ne0\) and its spectrum is different from \((a,a,-a,-a)\)
up to ordering, at most one of \(p,q,r\) is zero. The trace
identity for \(w_{ij}\) then makes every sectional Weyl coefficient
zero in every Ricci eigenframe.

The condition \(C=0\) says that the tangent-valued one-form
\(Z^\sharp\) is covariantly closed. Apply the covariant exterior
derivative again and use a Ricci eigenframe. With the curvature
indices ordered by \(R_{ijkl}=\langle R(e_i,e_j)e_l,e_k\rangle\),
the resulting identity is
\[
 \xi_kR_{ijkl}+\xi_iR_{jkil}+\xi_jR_{kijl}=0 .
\]
For \(l=i\), ordinary Bianchi gives
\((\xi_k-\xi_j)R_{ijki}=0\). Thus a Weyl coefficient with exactly
one shared index is zero when the other two eigenvalues differ,
since \(\Rm-W\) is diagonal in the Ricci eigenframe. When they
coincide, polarize the zero sectional Weyl coefficient using a
rotation in their common eigenspace. For entries with four
distinct indices, set \(x=R_{1234}\), \(y=R_{2314}\), and
\(z=R_{3124}\). Bianchi gives \(x+y+z=0\). The covariant closure
identity with last indices 4 and 3 gives, respectively,
\(\xi_3x+\xi_1y+\xi_2z=0\) and
\(\xi_4x+\xi_2y+\xi_1z=0\). Adding gives \(px=0\).
Permuting indices gives \(qz=ry=0\), so at most one of \(x,y,z\)
can be nonzero and Bianchi forces all three to vanish. Thus
\(W=0\) at every point where \(Z\ne0\) has spectrum different
from the double pairs.

The open set on which both \(W\) and \(Z\) are nonzero therefore
has Ricci defect eigenvalues \(a,-a\) of multiplicity two, with
\(a>0\) smooth. For either repeated eigenvalue, take two
orthogonal unit vectors in its eigenspace. Pair the Codazzi
equation with one of them to see that the eigenvalue derivative
along the other is zero. Applying this in the two eigenspaces
makes \(a\) locally constant. For inputs from opposite
eigenspaces the two sides of the Codazzi equation lie in
orthogonal summands and must both vanish. This gives parallel
eigenspace distributions and hence a local product of surfaces.
Its Gaussian curvatures are \(R/4+a\) and \(R/4-a\). If the first
two indices span the first factor, the decomposition of curvature
gives \(w_{12}=R/6\) and \(p=2a\), contradicting \(p w_{12}=0\).
Consequently wherever \(Z\ne0\), one has \(W=0\).

If \(Z\) is nonzero anywhere, \(W\) thus vanishes on an open set.
The unique continuation theorem for local conformal flatness of
smooth connected obstruction-flat metrics applies, since
obstruction-flatness is Bach-flatness in four dimensions
\cite[Theorem~3.8]{LassasLiimatainen2019ConformalHarmonic}.
It gives \(W=0\) globally. On a closed simply connected four-sphere
type, Kuiper's theorem supplies a conformal diffeomorphism to the
round sphere \cite{Kuiper1949}. Write the pulled-back round metric
as \(f^{-2}g\) with \(f>0\). The trace-free Ricci transformation
then gives \(f Z=-2(\nabla^2 f)^0\). Pairing with \(Z\) and
integrating using \(\operatorname{div}Z=0\) shows that
\(\int f|Z|^2dV_g=0\). Hence \(Z=0\) and the positive conformally
flat metric has constant sectional curvature.
\end{proof}

The Cotton-flat subclass on the homotopy sphere accordingly
has zero in \eqref{eq:ricci-cotton-hess} and enters the Einstein
comparison below. Cotton flatness remains a separate hypothesis
from Bach stationarity and second-order stability.
